% 00-style-sample.tex – TEMPORARY design sample (Steps 09–10). Placeholder text only; not book content.
% Uses template commands only (no formatting defined here).

\ChapterMission{המשימה: להראות שהבעיה אמיתית, מדידה ומתרחשת כל בוקר}
\ChapterIcon{pedestrian}
\chapter{דוגמת סגנון: פרק בשלב הבעיה}
\begin{scenebox}
שבע וחצי בבוקר. האור האדום דולק, והכביש ריק לגמרי. אישה אחת עומדת לבדה על המדרכה,
מביטה ימינה ושמאלה, ומעבירה את משקלה מרגל לרגל. אין איש לידה שיראה אותה.
היא יורדת מהמדרכה.
\end{scenebox}

זהו טקסט לדוגמה שנועד לבדוק את עיצוב העמוד בלבד, ואינו חלק מתוכן הדוח.\sidenote{הערת שוליים לדוגמה: הערות קצרות מופיעות בשוליים החיצוניים.}
הפסקה משלבת עברית עם מונחים באנגלית, כמו \en{Smart Pedestrian Crossing} ו-\en{countdown signal},
וכן מספרים: \pct{13.5} מההולכים שהגיעו במופע האדום חצו באדום, מתוך \num{563} תצפיות במופע זה
ו-\num{1,392} תצפיות בסך הכול \cite{rosenbloom2009}. רווח הסמך לדוגמה הוא \range{0.51}{0.78},
וזמן ההמתנה המרבי המתוכנן הוא \unit{30}{s}.

\begin{statrow}[3]
  \stattile[SignalRed]{\pct{13.5}}{מהמגיעים במופע האדום חצו באדום}
  \stattile[SignalAmber]{\num{1,392}}{הולכי רגל נצפו בצומת אחד בשעת בוקר}
  \stattile[SignalGreen]{\num{0.63}}{יחס הסיכויים לכל ממתין נוסף (מחושב מחדש)}
\end{statrow}

\begin{pullquote}[ציטוט לדוגמה, לא מתוך המקור]
כשאין איש על המדרכה, הרמזור הוא רק הצעה.
\end{pullquote}

\section{סעיף לדוגמה עם תיבות}
משפט מעבר קצר מציג את התיבות. כל תיבה ממלאת תפקיד קבוע לאורך הדוח, וצבעה נגזר מפלטת הרמזור.

\begin{findingbox}
ממצא לדוגמה: ככל שממתינים יותר אנשים על המדרכה, כך יורד הסיכוי לחצייה באדום.
\end{findingbox}

\begin{cautionbox}
סוגיה לדוגמה: עמודת \en{$e^{b}$} בטבלה המקורית אינה תואמת את \en{$\exp(b)$} בשני משתנים.
\end{cautionbox}

\begin{trapbox}
מלכודת לדוגמה: ממחקר תצפיתי אין להסיק קשר סיבתי.
\end{trapbox}

\begin{definitionbox}[הגדרה: שליטה חברתית]
הגדרה לדוגמה: מנגנון שבו נוכחות של אחרים מעודדת את הפרט לפעול לפי הנורמה.
\end{definitionbox}

\begin{projectbox}
יישום לדוגמה: חיישן המדרכה סופר את הממתינים ומעביר את הנתון לבקר הרמזור.
\end{projectbox}

\subsection{תת-סעיף לדוגמה}
רשימה לדוגמה:
\begin{itemize}
  \item פריט ראשון עם מספר: \pct{40.4} הגיעו במופע האדום.
  \item פריט שני עם מונח באנגלית: \en{Hosmer–Lemeshow}.
  \item פריט שלישי עם גרסה: \en{LuaHBTeX} \num{1.24.0}.
\end{itemize}
\crosswalkrule

\setcounter{chapter}{3}
\ChapterMission{המשימה: לבחון את מאמר הדגל בעין ביקורתית ולזהות את פער הידע}
\ChapterIcon{sensor}
\chapter{דוגמת סגנון: פרק בשלב הניתוח}
\begin{scenebox}
שמונה בבוקר, אותו מעבר חציה. על המדרכה עומדים כבר שבעה אנשים. אף אחד לא זז.
מישהו בודק את הטלפון, מישהו אחר מביט ברמזור. האור מתחלף, וכולם חוצים יחד.
\end{scenebox}

פסקה לדוגמה מפנה לטבלה~\ref{tab:sample}, למשוואה~\eqref{eq:logit} ולאיור~\ref{fig:photo}. הטבלה
מציגה נתונים שעובדו מן המקור \cite{rosenbloom2009}, והמשוואה מציגה את מבנה המודל הלוגיסטי.%
\sidenote{המספרים בטבלה חושבו מחדש מטבלה~1 במקור.}

\begin{table}[ht]
  \centering
  \caption{טבלה לדוגמה: חצייה באדום בקרב המגיעים במופע האדום, לפי מגדר (עובד מתוך \cite{rosenbloom2009})}
  \label{tab:sample}
  \begin{tabular}{lccc}
    \toprule
    \headrow \thead{מגדר} & \thead{הגיעו באדום} & \thead{חצו באדום} & \thead{שיעור} \\
    \midrule
    נשים & \num{330} & \num{28} & \pct{8.5} \\
    גברים & \num{233} & \num{48} & \pct{20.6} \\
    סך הכול & \num{563} & \num{76} & \pct{13.5} \\
    \bottomrule
  \end{tabular}
\end{table}

\begin{equation}
  \ln\frac{p}{1-p} = \beta_0 + \beta_1 x_{\hebmath{מגדר}} + \beta_2 x_{\hebmath{ממתינים}} + \dots + \beta_k x_k
  \label{eq:logit}
\end{equation}

במשוואה~\eqref{eq:logit}, \en{$p$} הוא ההסתברות לחצייה באדום, ויחס הסיכויים של כל משתנה הוא
\en{$\mathrm{OR}=e^{\beta}$}. לדוגמה, עבור \en{$b=-0.47$} מתקבל \en{$e^{b}\approx 0.63$}.

\authorphoto[48mm]{figures/photo-crosswalk.jpg}{מעבר החציה שנבחר לפרויקט, בשעת בוקר}{fig:photo}
\crosswalkrule
